\documentclass{article}
\title{線形代数 応用編 — 固有値から主成分分析へ}
\author{TeX 図書館}

\begin{document}

\section{この教科書のために — 固有値は実データで働く}

『大学数学入門』で学んだ固有値・固有ベクトルは、計算の練習題材のように見えるかもしれません。この応用編では、\textbf{固有値が実データの分析エンジン}として働く姿を見ます。到達点は\textbf{主成分分析}( PCA )と\textbf{特異値分解}( SVD )——データ科学の屋台骨です。

\section{対称行列の幸運 — 直交対角化}

実\textbf{対称}行列(\(A = A^{\top}\))には、素晴らしい性質があります。

\begin{quote}
\textbf{【定理:直交対角化】}\ 実対称行列は、直交行列 \(Q\)(\(Q^{\top}Q = I\))を使って
\[ A = Q D Q^{\top} \]
と対角化できる。\(D\) は実固有値を対角に並べたもので、異なる固有値に対応する固有ベクトルは\textbf{互いに直交}する。
\end{quote}

一般の行列は対角化できないこともありましたが、\textbf{対称行列なら常に}です。しかも固有ベクトルが直交——「互いに独立な軸」が自動で手に入ります。データ分析で共分散行列を扱うとき、この性質が全部を支えます。

\begin{quote}
\textbf{【演習】}\ 対称行列 \(A = \begin{pmatrix} 2 & 1 \\ 1 & 2 \end{pmatrix}\) を直交対角化せよ(固有値 1, 3)。
\end{quote}

\textbf{解答の指針:}\ 固有ベクトルは \(\begin{pmatrix}1\\-1\end{pmatrix}, \begin{pmatrix}1\\1\end{pmatrix}\)、長さで割って直交行列 \(Q = \frac{1}{\sqrt{2}}\begin{pmatrix} 1 & 1 \\ -1 & 1 \end{pmatrix}\)。\(Q^{\top}AQ = \mathrm{diag}(1, 3)\)。

\section{二次形式と等高線 — 楕円の向きを決めるもの}

対称行列 \(A\) とベクトル \(\boldsymbol{x}\) の組
\[ f(\boldsymbol{x}) = \boldsymbol{x}^{\top} A \boldsymbol{x} \]
を\textbf{二次形式}といいます。2 変数なら \(f = ax^2 + 2bxy + cy^2\) で、その等高線は楕円か双曲線です。

直交対角化すると見通しが劇的に良くなります:\(\boldsymbol{x} = Q\boldsymbol{y}\) と座標変換すると
\[ f = \boldsymbol{y}^{\top} D \boldsymbol{y} = \lambda_1 y_1^2 + \lambda_2 y_2^2 \]
——\textbf{交差項が消え、軸が固有ベクトルの方向に揃う}。つまり\textbf{固有ベクトルは楕円の主軸の向き}、\textbf{固有値は軸の長さの指標}です。「データの楕円」の向きと伸び方を固有値が語る、が PCA の幾何学的イメージです。

\begin{quote}
\textbf{【演習】}\ \(A = \begin{pmatrix} 5 & 2 \\ 2 & 5 \end{pmatrix}\) の二次形式の等高線の主軸の向きを固有ベクトルで答えよ。
\end{quote}

\textbf{解答の指針:}\ 固有値 7 と 3、固有ベクトル \((1,1)/\sqrt{2}\) と \((1,-1)/\sqrt{2}\)。長軸(固有値小=縮みにくい方向の解釈にも注意)は \((1,1)\) 方向と \((1,-1)\) 方向に揃う。

\section{主成分分析 — 分散が最大の方向}

データが \(n\) 点あるとします(各点は \(p\) 次元ベクトル)。\textbf{主成分分析}( PCA )の問いは:

\begin{quote}
「データを 1 本の直線に要約するなら、\textbf{どの方向で射影するのが最も情報を保つか}?」
\end{quote}

答えは\textbf{射影後の分散が最大の方向}です。そしてこの方向は、\textbf{平均を引いたデータの共分散行列 \(S\) の最大固有値の固有ベクトル}に一致します。

なぜそうなるかの骨子だけ示します。方向 \(\boldsymbol{u}\)(単位ベクトル)への射影の分散は
\[ \mathrm{Var}(\boldsymbol{u}^{\top}\boldsymbol{x}) = \boldsymbol{u}^{\top} S \boldsymbol{u} \]
これを \(\|\boldsymbol{u}\| = 1\) の制約つきで最大化すると、ラグランジュの未定乗数法により
\[ S\boldsymbol{u} = \lambda \boldsymbol{u} \]
——\textbf{固有値問題そのもの}です。最大固有値の固有ベクトルが最適方向(第 1 主成分)、次の固有値が第 2 主成分(第 1 と直交)、…と続きます。

\begin{quote}
\textbf{【例】}\ 「身長・体重」のデータの共分散行列が \(S = \begin{pmatrix} 100 & 60 \\ 60 & 50 \end{pmatrix}\) だったとします。固有値・固有ベクトルを計算すると、最大固有値の固有ベクトルは「身長も体重も増える方向」。第 1 主成分は「\textbf{体格}」という要約軸、第 2 主成分(直交方向)は「身長の割に体重?」という\textbf{体型差}の軸になります。
\end{quote}

\begin{quote}
\textbf{【演習】}\ 上の \(S\) の固有値を求め、第 1 主成分の「説明分散の割合」を計算せよ(固有値の総和に対する最大固有値の割合)。
\end{quote}

\textbf{解答の指針:}\ 固有方程式 \((100-\lambda)(50-\lambda) - 3600 = 0\) は \(\lambda^2 - 150\lambda + 1400 = 0\)、判別式 \(150^2 - 5600 = 16900 = 130^2\) より \(\lambda = 140, 10\)。割合は \(140/150 \approx 93\%\)。\textbf{2 変数の情報の 93\% が 1 本の軸に圧縮できる}。

\section{PCA の実務の読み方 — 寄与率と次元削減}

\begin{enumerate}
\item \textbf{寄与率}: 各固有値 ÷ 固有値の総和。「その主成分が説明するばらつきの割合」。累積寄与率が 80〜90\% に達するまで主成分を採用するのが定番。
\item \textbf{次元削減}: \(p\) 変数を「上位 \(k\) 主成分」に置き換える。可視化(2〜3 次元へ)・ノイズ除去・後段モデルの安定化に効く。
\item \textbf{注意}: PCA は\textbf{尺度に敏感}。単位の違う変数(円と cm)が混ざると、数値が大きい変数が支配するので、事前に標準化(平均 0・分散 1)するのが標準です。
\end{enumerate}

\begin{quote}
\textbf{【演習】}\ 固有値 \(100, 50, 30, 20\) のとき、累積寄与率 90\% を超える最小の主成分数を求めよ。
\end{quote}

\textbf{解答の指針:}\ 総和 200。上位 3 個で \(180/200 = 90\%\)。よって 3 個。

\section{特異値分解 — どんな行列でも分解できる}

対称でない行列も分解したい。そこで\textbf{特異値分解}( SVD )です。任意の \(m \times n\) 行列は
\[ A = U \Sigma V^{\top} \]
と書けます。\(U, V\) は直交行列、\(\Sigma\) は「\textbf{特異値} \(\sigma_i \ge 0\) を大きい順に対角に置いた」行列です。

特異値は \(A^{\top}A\)(これは対称!)の固有値の\textbf{平方根}なので、対称行列の理論がそのまま効いています。

\begin{quote}
\textbf{【定理:特異値と作用の幾何】}\ 単位球を \(A\) で写すと、楕円体になる。その\textbf{主軸の向きが \(V\) の列(右特異ベクトル)}、\textbf{軸の長さが特異値}、写像先の向きが \(U\) の列。固有値の「楕円の主軸」の話が、任意行列に一般化された形です。
\end{quote}

\begin{quote}
\textbf{【演習】}\ \(A^{\top}A\) の固有値が \(25, 9, 1\) のとき、\(A\) の特異値をすべて答えよ。
\end{quote}

\textbf{解答の指針:}\ 平方根をとるので \(5, 3, 1\)。

\section{SVD の応用 — 圧縮と推薦}

SVD の威力は\textbf{低ランク近似}です。特異値は大きい順なので、上位 \(k\) 個だけ残した
\[ A_k = \sum_{i=1}^{k} \sigma_i \boldsymbol{u}_i \boldsymbol{v}_i^{\top} \]
が「\textbf{フロベニウスノルムの意味で最良の}階数 \(k\) 近似」になります(エッカート・ヤングの定理)。

\begin{itemize}
\item \textbf{画像圧縮}: 画像 = 巨大行列。上位特異値だけで再構成すると、見た目を保ったまま情報量が激減する(ノイズは小さい特異値に追いやられる)。
\item \textbf{推薦システム}: 「ユーザー × 商品の評価」行列を SVD すると、ユーザーと商品を\textbf{潜在因子}(ジャンル好みなど)の空間に置ける。欠けている評価の予測値 = 低ランク近似の成分。
\item \textbf{ノイズ除去}: データ行列の上位特異値 = 信号、下位 = ノイズ、という切り分け。
\end{itemize}

\begin{quote}
\textbf{【演習】}\ 特異値 \(100, 50, 10, 1\) の行列を上位 2 項で近似したとき、残したエネルギーの割合(特異値 2 乗の比)を求めよ。
\end{quote}

\textbf{解答の指針:}\ \((100^2 + 50^2)/(100^2 + 50^2 + 10^2 + 1) = 12500/12601 \approx 99.2\%\)。特異値の\textbf{2 乗}比で測るのがポイント。

\section{マルコフ連鎖の定常分布 — 固有値 1 の固有ベクトル}

『確率過程入門』で出会った定常分布 \(\pi P = \pi\) は、線形代数の言葉では「\textbf{固有値 1 の左固有ベクトル}」です。

遷移行列の固有値は |固有値| \(\le 1\) で、最大固有値は 1。定常分布は常に存在し、連結で非周期的なら「どの初期分布から始めても \(\boldsymbol{\pi}_0 P^n \to \pi\)」に収束します。\textbf{時間をかけると初期条件が忘れられる}——「マルコフ連鎖の長期は初期に依存しない」が固有ベクトルで保証されるのです。

\begin{quote}
\textbf{【演習】}\ \(P = \begin{pmatrix} 0.9 & 0.1 \\ 0.4 & 0.6 \end{pmatrix}\) の定常分布を求めよ。
\end{quote}

\textbf{解答の指針:}\ \(\pi P = \pi\) と \(\pi_1 + \pi_2 = 1\) から \(\pi = (0.8, 0.2)\)。直感的には「0.1 で出ていく速度」と「0.4 で戻ってくる速度」が釣り合う割合。

\section{おわりに — 分解の目}

この応用編で得た視点は「\textbf{複雑なデータも、良い軸に分解すれば単純になる}」です。

対称行列の直交対角化(§2)→ 二次形式の楕円の主軸(§3)→ PCA(§4・5)→ SVD(§6・7)→ 定常分布(§8)。固有値問題は「このデータ・この作用の\textbf{自然な座標軸}は何か」という一つの問いの、様々な着地でした。

次の一歩は、PCA の統計的側面(標準化・外れ値への頑健性)、SVD による文書検索(LSA)、そして『機械学習入門』の次元削減とモデル安定化です。

\begin{quote}
\textbf{【総合演習】}
\begin{enumerate}
\item 共分散行列 \(\begin{pmatrix} 4 & 2 \\ 2 & 4 \end{pmatrix}\) の主成分の方向と寄与率を求めよ。
\item 特異値 \(6, 3\) の行列のフロベニウスノルムを求めよ。
\item 「PCA の前に標準化すべき理由」を 1 文で述べよ。
\end{enumerate}
\end{quote}

\textbf{解答の指針:}\ 1. 固有値 6, 2、方向 \((1,1)/\sqrt{2}\) と \((1,-1)/\sqrt{2}\)、寄与率 \(75\% / 25\%\)。2. \(\sqrt{6^2 + 3^2} = \sqrt{45} = 3\sqrt{5}\)。3. 尺度の大きい変数が分散を支配し、真の構造を隠すから。

\end{document}
